\documentclass[10pt,a4paper,ragged2e,withhyper]{altacv}
%%%%%%%%%%%%%%%%%%%%%%%%%%%%%%%%%%%%%%%%%%%%%%%%%%%%%%%%%%%%%%%%%%%%%%%%%%%%%%%%%%%%%%%
% 0 -> Version orange
% 1 -> Version noir et blanc
% 2 -> Version verte et bleu
% 3 -> Version violette
% 5 -> Version polyFlag
% 6 -> Version rainbow
\newcommand{\valColour}{0}
\PassOptionsToPackage{dvipsnames}{xcolor}
\geometry{left=1.25cm,right=1.25cm,top=1.5cm,bottom=1.5cm,columnsep=1.2cm, 
% showframe
}
\usepackage{paracol}
\usepackage{ifthen}
\usepackage[french]{babel}
\usepackage{hyperref}

%%%%%%%%%%%%%%%%%%%%%%%%%%%%%%%%%%%%%%%%%%%%%%%%%%%%%%%%%%%%%%%%%%%%%%%%%%%%%%%%%%%%%%%
% Compilation
\ifxetexorluatex
  % If using xelatex or lualatex:
  \usepackage{fontspec}
  \setmainfont{Roboto Slab}
  \setsansfont{Lato}
  \renewcommand{\familydefault}{\sfdefault}
\else
  % If using pdflatex:
  \usepackage[utf8]{inputenc}
  \usepackage[T1]{fontenc}
	\usepackage[rm]{roboto}
  \usepackage[defaultsans]{lato}
  \renewcommand{\familydefault}{\sfdefault}
\fi

\usepackage{csquotes}

% Colours if you want to
\newcommand{\switchColour}[1]
{
  \ifthenelse{\equal{#1}{0}}{%ORANGE
    \definecolor{HeaderColor}{HTML}{FFA51F}
    \definecolor{ColumnColor}{HTML}{FFDEAD}
    \definecolor{NameColour}{HTML}{8A0808}
    \definecolor{TaglineColour}{HTML}{000000}
    \definecolor{HeadingColour}{HTML}{DF7401}
    \definecolor{HeadingRuleColour}{HTML}{DF7401}
    \definecolor{SubheadingColour}{HTML}{B40404}
    \definecolor{AccentColour}{HTML}{B40404}
    \definecolor{EmphasisColour}{HTML}{B40404}
    \definecolor{BodyColour}{HTML}{61210B}}{}
  \ifthenelse{\equal{#1}{1}}{%NOIR
    \definecolor{HeaderColor}{HTML}{919191}
    \definecolor{ColumnColor}{HTML}{D1D1D1}
    \definecolor{NameColour}{HTML}{000000}
    \definecolor{TaglineColour}{HTML}{000000}
    \definecolor{HeadingColour}{HTML}{000000}
    \definecolor{HeadingRuleColour}{HTML}{000000}
    \definecolor{SubheadingColour}{HTML}{000000}
    \definecolor{AccentColour}{HTML}{000000}
    \definecolor{EmphasisColour}{HTML}{000000}
    \definecolor{BodyColour}{HTML}{000000}}{}
  \ifthenelse{\equal{#1}{2}}{%VERTE
    \definecolor{HeaderColor}{HTML}{79DA7B}
    \definecolor{ColumnColor}{HTML}{34CCB3}
    \definecolor{NameColour}{HTML}{004D2E}
    \definecolor{TaglineColour}{HTML}{00663D}
    \definecolor{HeadingColour}{HTML}{00663D}
    \definecolor{HeadingRuleColour}{HTML}{00663D}
    \definecolor{SubheadingColour}{HTML}{00663D}
    \definecolor{AccentColour}{HTML}{00663D}
    \definecolor{EmphasisColour}{HTML}{00331E}
    \definecolor{BodyColour}{HTML}{001a0F}}{}
  \ifthenelse{\equal{#1}{3}}{%VIOLETTE
    \definecolor{HeaderColor}{HTML}{FFCCD8}
    \definecolor{ColumnColor}{HTML}{ECC2FF}
    \definecolor{NameColour}{HTML}{4D004D}%15
    \definecolor{TaglineColour}{HTML}{660066}%20
    \definecolor{HeadingColour}{HTML}{660066}
    \definecolor{HeadingRuleColour}{HTML}{660066}
    \definecolor{SubheadingColour}{HTML}{660066}
    \definecolor{AccentColour}{HTML}{660066}
    \definecolor{EmphasisColour}{HTML}{1a001a}%5
    \definecolor{BodyColour}{HTML}{330033}}{}%10
  \ifthenelse{\equal{#1}{4}}{%
    \definecolor{HeaderColor}{HTML}{357F2D}
    \definecolor{ColumnColor}{HTML}{B8E4B3}
    \definecolor{NameColour}{HTML}{000000}
    \definecolor{TaglineColour}{HTML}{000000}
    \definecolor{HeadingColour}{HTML}{000000}
    \definecolor{HeadingRuleColour}{HTML}{000000}
    \definecolor{SubheadingColour}{HTML}{000000}
    \definecolor{AccentColour}{HTML}{000000}
    \definecolor{EmphasisColour}{HTML}{000000}
    \definecolor{BodyColour}{HTML}{000000}}{}
  \ifthenelse{\equal{#1}{5}}{%POLYFLAG
    \definecolor{HeaderColor}{HTML}{357F2D}
    \definecolor{ColumnColor}{HTML}{B8E4B3}
    \definecolor{NameColour}{HTML}{DF0101}
    \definecolor{TaglineColour}{HTML}{FF0040}
    \definecolor{HeadingColour}{HTML}{0B0B61}
    \definecolor{HeadingRuleColour}{HTML}{610B4B}
    \definecolor{SubheadingColour}{HTML}{0B0B61}
    \definecolor{AccentColour}{HTML}{0000FF}
    \definecolor{EmphasisColour}{HTML}{0B0B61}
    \definecolor{BodyColour}{HTML}{0B2161}}{}
  \ifthenelse{\equal{#1}{6}}{%RAINBOW
    \definecolor{HeaderColor}{HTML}{357F2D}
    \definecolor{ColumnColor}{HTML}{B8E4B3}
    \definecolor{NameColour}{HTML}{000000}
    \definecolor{TaglineColour}{HTML}{000000}
    \definecolor{HeadingColour}{HTML}{000000}
    \definecolor{HeadingRuleColour}{HTML}{000000}
    \definecolor{SubheadingColour}{HTML}{000000}
    \definecolor{AccentColour}{HTML}{000000}
    \definecolor{EmphasisColour}{HTML}{000000}
    \definecolor{BodyColour}{HTML}{000000}}{}
}
\switchColour{\valColour}

\colorlet{name}{NameColour}
\colorlet{tagline}{TaglineColour}
\colorlet{heading}{HeadingColour}
\colorlet{headingrule}{HeadingRuleColour}
\colorlet{subheading}{SubheadingColour}
\colorlet{accent}{AccentColour}
\colorlet{emphasis}{EmphasisColour}
\colorlet{body}{BodyColour}

% Fonts
\renewcommand{\namefont}{\Huge\rmfamily\bfseries}
\renewcommand{\personalinfofont}{\footnotesize}
\renewcommand{\cvsectionfont}{\LARGE\rmfamily\bfseries}
\renewcommand{\cvsubsectionfont}{\large\bfseries}


% Bullets for itemize and rating marker
\renewcommand{\itemmarker}{{\small\textbullet}}
\renewcommand{\ratingmarker}{\faCircle}

\newcommand{\headerbg}[1]{
    \begin{tikzpicture}[remember picture,overlay]
      \fill[HeaderColor] (current page.north west) rectangle ++(\paperwidth,-#1);
    \end{tikzpicture}
}

\newcommand{\rightcolumnbg}[1]{
    \begin{tikzpicture}[remember picture,overlay]
      \fill[ColumnColor] (current page.north east) rectangle ++(-#1\paperwidth-0.05cm,-\paperheight);
    \end{tikzpicture}
}

\newcommand{\leftcolumnbg}[1]{
    \begin{tikzpicture}[remember picture,overlay]
      \fill[ColumnColor] (current page.north west) rectangle ++(#1\paperwidth+0.05cm,-\paperheight);
    \end{tikzpicture}
}


%%%%%%%%%%%%%%%%%%%%%%%%%%%%%%%%%%%%%%%%%%%%%%%%%%%%%%%%%%%%%%%%%%%%%%%%%%%%%%%%%%%%%%%
\begin{document}
\name{Prunelle Daudré--Treuil}
\tagline{Étudiante en Master TAL-NLP recherchant une thèse\hskip 1cm 26 ans} %/!\ À CHANGER
\photoR{2.8cm}{picme}

\personalinfo{
  \email{prunelle.daudretreuil@protonmail.com}
  \phone{+33(0)7 81 80 11 77}
  \homepage{ptreuil.github.io}
  \mailaddress{7 Rue de l'Armée Patton}
  \location{Nancy, FRANCE}
  \github{ptreuil}
  \linkedin{ptreuil}
  
}

%%%%%%%%%%%%%%%%%%%%%%%%%%%%%%%%%%%%%%%%%%%%%%%%%%%%%%%%%%%%%%%%%%%%%%%%%%%%%%%%%%%%%%%
\rightcolumnbg{0.4} % les commandes sont toujours à mettre dans ce sens ! les params de columnbg + paracol doivent donner 1
\headerbg{4.95cm}

\makecvheader

\columnratio{0.6}

\begin{paracol}{2}

\cvsection{Expériences professionnelles}

\cvevent{Doctorat en Informatique}{LORIA}{oct 2024 -- sept 2027}{Vandoeuvre-lès-Nancy (54), FRANCE}

Construction d'un système conversationnel pour un accès multi-modal à un corpus

\divider

\cvevent{Stage de fin d'étude (Master TAL)}{LORIA}{mars 2024 -- août 2024}{Vandoeuvre-lès-Nancy (54), FRANCE}

Exploitation et appréhension de la correspondance d'Henri Poincaré

% Sous la direction de M.Nauer, M.Lieber et M.D'Aquin du LORIA et M.Bruneau et M.Rollet des Archives Henri Poincaré. Réalisation d'un système d’interrogation conversationnelle permettant un accès hiérarchiques aux lettres du mathématicien Henri Poincaré en fonction des propriétés de celles-ci

\vspace{0.2cm}

Présentation orale à la \textcolor{AccentColour}{\href{https://afia.asso.fr/les-journees-communes/humanite-numerique-et-ia/}{Journée Humanités Numériques et IA}} organisée par l'AFIA et les GdR MADICS et MAGIS à la BNF le 3 mais 2024

\divider

\cvevent{Stage d'observation (Licence MIASHS)}{Université de Lille}{mai 2021 -- juin 2021}{Tourcoing (59), FRANCE}
Au sein du laboratoire SCALAB auprès de Mme.Delevoye, professeur en psychologie cognitive

\divider

\cvevent{Service Civique}{Unis-cité}{oct 2018 -- juin 2019}{Belfort (90), FRANCE}
Missions Cinéma et Citoyenneté \& Jeunes Citoyens du Numérique

% \vspace{0.2cm}

\cvsection{Projets}

% \cvprojecthub{Stage M2 : Conversational Recommender System}{}{Archives Poincaré}{mars 2024 -- août 2024}{Dévellopement d'un outil de fouille d'un corpus structuré en graphe de connaissance}

% \divider

\cvprojecthub{Projet M2 : Raisonnement à Partir de Cas par analogie}{}{Université de Lorraine}{2023}{Analyse des résultats d'un algorithme de RàPC utilisant l'apprentissage par analogie}

\divider

\cvprojecthub{Projet M1 : Analyse des variations sur les manuscrits  du livre de Ben-Sira}{}{Université de Lorraine}{2022-2023}{Analyse des différences de transcription entre deux manuscrits de la Bible en ancien hébreu}

\divider

\cvprojecthub{Projet M1 : Extraction d'information depuis des articles scientifiques}{}{Université de Lorraine}{Automne 2022}{}

\switchcolumn

\cvsection{Formations}

\cvachievement{\faUniversity}{Master TAL-NLP/Msc NLP}{Université de Lorraine \hspace{0.2cm}(2022--2024)}

\vspace{0.2cm}

\divider

\cvachievement{\faUniversity}{Licence MIASHS,  Parcours Sciences Cognitives}{Université de Lille \hspace{0.9cm}(2019-2022)}

\vspace{0.2cm}

\divider

\cvachievement{\faUniversity}{Formation d'ingénieur - Informatique}{UTBM \hspace{2.6cm}(2014-2018)}

\divider

\cvachievement{\faStar}{TOIEC (2024)}{CEFR Level : C1, Score : 980}

\cvsection{Langues}

\cvskill{Français}{5}

\divider

\cvskill{Anglais}{4}

\vspace{0.3cm}

\cvsection{Compétences}

\cvtag{Python}
\cvtag{Numpy}
\cvtag{Panda}
\cvtag{Matplotlib}
\cvtag{RdfLib}
\cvtag{SPARQLWrapper}
\cvtag{SKLearn}
\cvtag{FastText}
\cvtag{Spacy}
\cvtag{NLTK}
\cvtag{Stanza}
\cvtag{Pytorch}
\cvtag{Networkx}

\divider

\cvtag{RDF}
\cvtag{RDFS}
\cvtag{OWL}
\cvtag{SPARQL}\\
\cvtag{Protégé}

\divider

\cvtag{\LaTeX}
\cvtag{Zotero}
\cvtag{Bash}
\cvtag{Git}
\cvtag{HTML}
\cvtag{Javascript}
\vspace{0.3cm}

\cvsection{Centres d'intérêt}

\hspace{-1.5cm}
\wheelchart{1.5cm}{0.5cm}{
 4/7em/accent!45/{Pyrogravure},
 4/7em/accent!50/{Joaillerie},
 4/7em/accent!55/{Jeux vidéos},
 4/7em/accent!60/{Littérature},
 4/7em/accent!65/{Marche},
 4/7em/accent!70/{\LaTeX},
 4/7em/accent!75/{Voyage}
}

\end{paracol}

\end{document}
